\documentclass[12pt,a4paper,UTF8]{ctexart}
\usepackage[a4paper,top=22mm,bottom=22mm,left=22mm,right=22mm,headheight=15pt]{geometry}
\newcommand{\Needspace}[1]{\par\begingroup\dimen0=#1\relax\dimen2=\pagegoal\advance\dimen2 by-\pagetotal\ifdim\dimen0>\dimen2\newpage\fi\endgroup}
\newcommand{\needspace}[1]{\par\begingroup\skip0=\lastskip\vskip-\skip0\dimen0=#1\relax\advance\dimen0 by\skip0\vskip\dimen0\penalty-100\vskip-\dimen0\vskip\skip0\penalty9999\endgroup}
\usepackage{amsmath,amssymb,mathtools,bm}
\usepackage{graphicx,xcolor,booktabs,longtable,array,enumitem,fancyhdr,caption,float,microtype,setspace}
\usepackage[hidelinks,unicode]{hyperref}
\usepackage[most]{tcolorbox}
\definecolor{ContestBlue}{HTML}{17365D}
\hypersetup{pdftitle={第074套},pdfauthor={}}
\setstretch{1.24}
\setlength{\parindent}{2em}
\setlength{\parskip}{0.35em}
\setlength{\emergencystretch}{2em}
\setlist{itemsep=0.35em,topsep=0.4em}
\captionsetup{font=small,labelfont=bf,justification=centering}
\setcounter{secnumdepth}{0}
\pagestyle{fancy}
\fancyhf{}
\fancyhead[L]{\small 第074套}
\fancyfoot[C]{\thepage}
\renewcommand{\headrulewidth}{0.4pt}
\ctexset{
  section={format=\Large\bfseries\color{ContestBlue},beforeskip=1.2em,afterskip=0.7em},
  subsection={format=\large\bfseries\color{ContestBlue},beforeskip=1.0em,afterskip=0.5em},
  subsubsection={format=\normalsize\bfseries,beforeskip=0.8em,afterskip=0.35em}
}
\usepackage{calc}
\newcounter{none}
\let\originalsubsection\subsection
\renewcommand{\subsection}{\needspace{6\baselineskip}\originalsubsection}
\makeatletter
\newcommand{\tightlist}{\setlength{\itemsep}{0pt}\setlength{\parskip}{0pt}}
\newcommand{\maxwidth}{\ifdim\Gin@nat@width>\linewidth\linewidth\else\Gin@nat@width\fi}
\newcommand{\maxheight}{\ifdim\Gin@nat@height>0.72\textheight 0.72\textheight\else\Gin@nat@height\fi}
\newcommand{\pandocbounded}[1]{\begingroup\setkeys{Gin}{width=\maxwidth,height=\maxheight,keepaspectratio}#1\endgroup}
\makeatother
\begin{document}
\begin{titlepage}
  \centering
  \vspace*{58mm}
  {\Huge\bfseries\color{ContestBlue} 第074套\par}
  \vfill
  \vspace*{28mm}
\end{titlepage}
\section{第 1 题}\label{ux7b2c-1-ux9898}

一根长度为 \(L=1\ \mathrm m\)
的弹性轻杆水平悬臂放置，左端固定。在右端悬挂质量为 \(1\ \mathrm{kg}\)
的重物后，杆的右端下降
\(h=1\ \mathrm{cm}\)。现将这根杆竖直放在地面上，下端固定，在上端施加竖直向下的压力
\(F\)。估算使杆开始发生弯曲失稳所需的力 \(F\)。

\begin{figure}[H]
\centering
\includegraphics[width=0.322\linewidth,height=\textheight,keepaspectratio,alt={图1-1　竖直弹性杆的受压失稳模型}]{../figures/question-01.pdf}
\caption*{图1-1　竖直弹性杆的受压失稳模型}
\end{figure}

\newpage

\section{第 2 题}\label{ux7b2c-2-ux9898}

在光滑水平面上有三个半径均为 \(R\) 的均匀球。球1和球2的质量均为
\(m\)，球3的质量为
\(\lambda m\)。初始时球2静止叠放在球1正上方，球3以速度 \(v\)
向右运动。球与球之间的碰撞均为完全弹性碰撞，所有接触面均无摩擦。

（1）求球3与球1第一次碰撞结束瞬间两球的速度，以向右为正方向。

（2）暂不考虑球2与球3的碰撞。求 \(\lambda\) 与 \(v\)
满足什么条件时，第一次碰撞后球2会直接进入自由落体，而不是沿球1表面滑下。

（3）在满足第（2）问的条件下，球2随后与球3相碰，且碰撞瞬间两球球心连线与水平方向成
\(45^\circ\)。求这种情形可能对应的 \(\lambda\) 范围。

（4）取
\(\lambda=3\)、\(v=\sqrt{Rg}\)，并假设球3始终不离开地面。求球2此后相对地面所能达到的最大高度
\(h\)，以 \(R\) 表示。

\begin{figure}[H]
\centering
\includegraphics[width=0.631\linewidth,height=\textheight,keepaspectratio,alt={图2-1　三球的初始位置与第一次碰撞}]{../figures/question-02.pdf}
\caption*{图2-1　三球的初始位置与第一次碰撞}
\end{figure}

\newpage

\section{第 3 题}\label{ux7b2c-3-ux9898}

两根相距 \(l\) 的无限长带电直线均平行于 \(z\)
轴固定放置，线电荷密度分别为 \(+\lambda\) 和 \(-\lambda\)。在
\(x\)---\(y\)
平面内，负线电荷位于左侧，正线电荷位于右侧，连接两线电荷的主轴正方向向右。取两线电荷中点为极坐标原点。

（1）当 \(r\gg l\) 时，求点 \(A(r,\theta)\) 的电势。

（2）求该点的径向电场分量 \(E_r\) 和角向电场分量
\(E_\theta\)，并说明等势线与电场线的形状。

（3）移去线电荷 \(-\lambda\)，并将线电荷 \(+\lambda\) 平移到原点。在
\((0,r_0)\) 处放置质量为 \(m\)、电荷量为 \(-q\) 的粒子，给它沿 \(+x\)
方向的初速度

\[
v_0=\sqrt{\frac{2\lambda q}{\pi\varepsilon_0m}}.
\]

求：

I. 粒子的径向速度下一次等于零时，它到原点的距离 \(r_1\)；

\begin{enumerate}
\def\labelenumi{\Roman{enumi}.}
\setcounter{enumi}{1}
\tightlist
\item
  从初始时刻到该时刻，粒子绕原点转过的角度 \(\theta_0\)。
\end{enumerate}

\begin{figure}[H]
\centering
\includegraphics[width=0.628\linewidth,height=\textheight,keepaspectratio,alt={图3-1　异号线电荷在横截面内的位置}]{../figures/question-03.pdf}
\caption*{图3-1　异号线电荷在横截面内的位置}
\end{figure}

\newpage

\section{第 4 题}\label{ux7b2c-4-ux9898}

边长为 \(L\) 的正方形平行板电容器，两极板间距为
\(d\)。先把电容器接到电压为 \(V\)
的电池上，充电完成后断开电池；再将电容器竖直插入相对介电常数为
\(\varepsilon_r\)、密度为 \(\rho\)
的液体中。稳定时，液体在两极板间所占的竖直长度为
\(L/2\)。忽略其他外力影响，求：

（1）此时电容器的电容；

（2）两极板间的电场强度；

（3）金属极板自由电荷的面密度；

（4）极板内外液面的高度差 \(h\)。

\begin{figure}[H]
\centering
\includegraphics[width=0.617\linewidth,height=\textheight,keepaspectratio,alt={图4-1　方形平行板电容器的竖直浸液状态}]{../figures/question-04.pdf}
\caption*{图4-1　方形平行板电容器的竖直浸液状态}
\end{figure}

\newpage

\section{第 5 题}\label{ux7b2c-5-ux9898}

半径为 \(R\)、质量为 \(M\) 的均匀圆盘起初静止在光滑水平桌面上。质量为
\(m\)、初速度为 \(\bm v_0\) 的质点射向圆盘，击中后嵌入盘缘上的 \(A\)
点。图中 \(C\) 为圆心，直线 \(C'C\) 平行于 \(\bm v_0\)，入射直线与
\(C'C\) 的距离为 \(b\)。点 \(B\) 位于盘缘，且 \(BC\perp C'C\)。已知
\(m\ll M\)。

（1）求使碰后 \(B\) 点相对桌面瞬时静止的 \(b\)。

（2）承接第（1）问，取质点嵌入的时刻为 \(t=0\)，以该时刻 \(B\)
点的位置为固定在桌面上的坐标原点，水平向右为 \(+x\)、竖直向上为
\(+y\)，求 \(A\) 点在时刻 \(t\) 的坐标 \(x(t)\)、\(y(t)\)。

（3）若在盘缘极靠近 \(B\)
点处开孔并以光滑小钉固定，使圆盘可绕小钉转动。嵌入过程中小钉一般承受冲量。求小钉冲量为零时的
\(b\)，并求此时圆盘的角速度。

已知均匀圆盘绕中心轴的转动惯量为 \(I_C=MR^2/2\)。

\begin{figure}[H]
\centering
\includegraphics[width=0.408\linewidth,height=\textheight,keepaspectratio,alt={图5-1　质点入射圆盘的偏心几何}]{../figures/question-05.pdf}
\caption*{图5-1　质点入射圆盘的偏心几何}
\end{figure}

\newpage

\section{第 6 题}\label{ux7b2c-6-ux9898}

一套蚊虫光学探测装置由薄凸透镜和位于其右侧的 CCD
探测器组成。蚊虫发出的特定频段射线经透镜折射后，只要有光线到达
CCD，即认为蚊虫被发现。取透镜平面中心为原点，光轴左方为
\(z>0\)，离轴距离记为 \(R=\sqrt{x^2+y^2}\)。透镜口径为
\(D=10\ \mathrm{cm}\)，焦距为 \(f=10\ \mathrm{cm}\)；CCD 与透镜相距
\(L=15\ \mathrm{cm}\)，CCD 直径为 \(d=4\ \mathrm{cm}\)。

（1）定量描述 \(z>0\) 区域内蚊虫不会被探测到的安全飞行区域，并在
\(R\)---\(z\) 截面内写出边界直线方程。

（2）蚊虫在不被发现的条件下，最小能接近到离光轴多远？

\begin{figure}[H]
\centering
\includegraphics[width=0.82\linewidth,height=\textheight,keepaspectratio,alt={图6-1　透镜、CCD 与空间点的光路几何}]{../figures/question-06.pdf}
\caption*{图6-1　透镜、CCD 与空间点的光路几何}
\end{figure}

\newpage

\section{第 7 题}\label{ux7b2c-7-ux9898}

一个内壁光滑的环形绝热气缸，其截面半径为 \(a\)，中心圆半径为
\(R\)。气缸内等间距放置 \(N\) 个质量均为 \(m\)
的绝热活塞，活塞与内壁无摩擦且不漏气。已知

\[
R\gg\frac RN\gg a,
\]

因此每两个相邻活塞之间的气室可近似为直圆柱。初态每个气室中有
\(\nu\ \mathrm{mol}\) 氮气，温度均为 \(T_0\)。记第 \(i\)
个活塞相对平衡位置的位移为 \(x_i\)，并有 \(|x_i|\ll R/N\)。

（1）求第1、2号活塞之间气体压强相对初态的变化量。

（2）考虑所有活塞都作角频率为 \(\omega\)
的简谐振动、相位随编号线性变化的基频行波模式。相邻活塞相位差
\(\Delta\varphi\ll1\)，可用
\(\cos\Delta\varphi\simeq1-(\Delta\varphi)^2/2\)。又因波速远小于气体声速，每个气室内压强可视为瞬时均匀。求波速的最小可能值。

\begin{figure}[H]
\centering
\includegraphics[width=0.731\linewidth,height=\textheight,keepaspectratio,alt={图7-1　环形绝热气缸与等间距活塞}]{../figures/question-07.pdf}
\caption*{图7-1　环形绝热气缸与等间距活塞}
\end{figure}

\newpage

\section{第 8 题}\label{ux7b2c-8-ux9898}

一束缪子在均匀磁场中沿圆形轨道运动，忽略电磁辐射造成的能量损失。缪子质量为

\[
m_\mu=1.88\times10^{-28}\ \mathrm{kg},
\]

电荷量为

\[
q_\mu=-1.602\times10^{-19}\ \mathrm C,
\]

半衰期为

\[
t_{1/2}=1.523\ \mu\mathrm s.
\]

（1）当缪子速度远小于光速时，若每运动一周恰有一半缪子衰变，求磁感应强度
\(B\)。

（2）保持同一磁场而提高缪子速度，使其不再远小于光速。问缪子每运动一周的衰变比例增大、减小还是保持不变？

\begin{figure}[H]
\centering
\includegraphics[width=0.515\linewidth,height=\textheight,keepaspectratio,alt={图8-1　缪子在均匀磁场中的圆周运动}]{../figures/question-08.pdf}
\caption*{图8-1　缪子在均匀磁场中的圆周运动}
\end{figure}
\end{document}
